\documentclass[10pt,a4paper]{amsart}
\usepackage[margin=19mm]{geometry}
\usepackage{amsmath,amssymb,mathtools}
\usepackage[scr=rsfs,cal=euler]{mathalfa}
\usepackage{xfrac}
\usepackage[unicode,hidelinks,bookmarks=false]{hyperref}
\newcommand{\ie}{i.e.\ }
\newcommand{\cf}{cf.\ }
\newcommand{\resp}{respectively\ }
\newcommand{\Subsection}{\subsection}
\providecommand{\nobreakdash}{\nobreak\hbox{-}}
\setlength{\emergencystretch}{2em}
\multlinegap0pt
\usepackage{amssymb}
\usepackage{xcolor}
\newtheorem{theorem}{Theorem}
\newtheorem{lemma}{Lemma}
\usepackage{cleveref}
\crefname{theorem}{Theorem}{Theorems}
\crefname{lemma}{Lemma}{Lemmas}
\newcommand{\R}{\mathbb R}
\DeclareMathOperator{\conv}{conv}
\title{Sequential testing of\\ conditionally constrained hypotheses}
\author{Eugenio Clerico}
\date{Source: arXiv:2606.06769v1 (CC BY 4.0)}
\begin{document}
\maketitle

\noindent\emph{Source and licence.} Sequential testing of\\ conditionally constrained hypotheses. Version arXiv:2606.06769v1 (CC BY 4.0), distributed by arXiv under the Creative Commons Attribution 4.0 International licence, http://creativecommons.org/licenses/by/4.0/, as recorded in the arXiv licence metadata for this exact version and shown on the abstract page https://arxiv.org/abs/2606.06769v1. Full licence notice: \texttt{licenses/arxiv-2606.06769-CC-BY-4.0.txt}.\par\smallskip

\noindent\textit{Editorial scope.} Counted unit: the article's lemma lemma:caradown (source0.tex lines 353--355). The article's complete original proof of it is delivered with it (source0.tex lines 356--376, one \texttt{proof} environment of 21 lines). No mathematics is written, repaired or re-derived: the statement and the proof are the article's own lines, in the article's own notation, and no retained line is substituted or re-flowed.  Retained uncounted as the source-local context the proof needs: lines 316--318.  Nothing was omitted from any retained span, so the package carries no omission record.  No retained line contains a citation, so the wrapper carries no bibliography and no bibliography entry.  No retained line contains a figure environment, an image-inclusion command or drawing code, so no image is delivered and the package declares no asset dependency.  Editorial formatting only, in addition: an AMS \texttt{amsart} preamble that restates the article's own declarations and macros used by the retained lines, namely the environments \texttt{theorem}, \texttt{lemma} on separate counters, together with the standard packages those lines need.  The retained environments print in this document as Theorem 1, Lemma 1 in order of appearance, whereas the article prints the same items with the numbering of its own full document.  The complete native source of the article as distributed in the arXiv e-print for this exact version is bundled unchanged as \texttt{sources/202211/source0.tex}.
\begin{theorem}[Carathéodory]\label{thm:caratheodory}
    Let $S\subseteq\R^m$. Then every point of $\conv S$ can be written as a convex combination of at most $m+1$ points of $S$, that is, as $\sum_{i=1}^k \alpha_i x_i$ for some $k\leq m+1$, $x_1,\dots,x_k\in S$, $\alpha_1,\dots,\alpha_k\geq 0$, and $\sum_{i=1}^k \alpha_i=1$.
\end{theorem}

\begin{lemma}\label{lemma:caradown}
    Let $A\subseteq\R^m\times\R$ be a non-empty set, and assume that whenever $(y,z)\in A$ and $z'\leq z$, then $(y,z')\in A$. Then, every point of $\conv A$ can be written as a convex combination of at most $m+1$ points of $A$.
\end{lemma}

\begin{proof}
    Let $(y,z)\in\conv A$. By \Cref{thm:caratheodory} in $\R^{m+1}$, we can write $(y,z)$ as a convex combination of at most $m+2$ points of $A$. If it requires $m+1$ or fewer points, we are done. 
    
    Thus, suppose $(y,z)=\sum_{i=1}^{m+2}(\alpha_i y_i,\alpha_i z_i)$ for points $(y_i,z_i)\in A$ and weights $\alpha_i>0$ summing to $1$. The $m+2$ points $y_1,\dots,y_{m+2}$ lie in $\R^m$, so they must be affinely dependent. Therefore, there exist real numbers $\gamma_1,\dots,\gamma_{m+2}$, not all zero, such that
    $$\sum_{i=1}^{m+2}\gamma_i y_i = 0 \qquad \text{and} \qquad \sum_{i=1}^{m+2}\gamma_i = 0\,.$$
    By replacing $\gamma$ with $-\gamma$ if necessary, we may assume without loss of generality that $\sum_{i=1}^{m+2}\gamma_i z_i \geq 0$. 
    
    Since $\sum_{i=1}^{m+2}\gamma_i=0$ and the $\gamma_i$ are not all zero, at least one $\gamma_i$ is strictly negative. Define
    $$t = \min_{\gamma_i < 0} \frac{\alpha_i}{-\gamma_i}\,.$$
    Because all $\alpha_i>0$, we have $t>0$. Let $j$ be an index where this minimum is attained. We define new weights $\beta_i = \alpha_i + t\gamma_i$. By our choice of $t$, we have $\beta_i \geq 0$ for all $i$, $\beta_j = 0$, and $\sum_{i=1}^{m+2}\beta_i = 1$. 
    
    Using these new weights, we have $\sum_{i=1}^{m+2}\beta_i y_i = y $. Let
    $$z^\star = \sum_{i=1}^{m+2}\beta_i z_i = z + t\sum_{i=1}^{m+2}\gamma_i z_i\,.$$
    Since $t>0$ and $\sum_{i=1}^{m+2}\gamma_i z_i \geq 0$, it follows that $z^\star \geq z$. 
    
    If $z^\star=z$, then $\sum_{i=1}^{m+2}\beta_i(y_i,z_i)$ is a convex combination of $(y,z)$ using at most $m+1$ points (since $\beta_j=0$), and we are done.
    
    If $z^\star>z$, pick any index $k$ such that $\beta_k>0$, and define $z_k' = z_k - \frac{z^\star-z}{\beta_k}$. Since $z_k'<z_k$, we have $(y_k,z_k')\in A$ by the downward closure assumption. Replacing $(y_k,z_k)$ with $(y_k,z_k')$ shifts the $z$-coordinate of the combination exactly to $z$:
    $$\sum_{i\neq k}\beta_i z_i + \beta_k z_k' = z^\star - (z^\star-z) = z\,.$$
    This yields $(y,z)$ as a convex combination of at most $m+1$ points of $A$.
\end{proof}

\end{document}
